\documentclass[12pt]{article}

\usepackage[a4paper, total={6in, 8in}]{geometry}
\usepackage{textcomp}

\begin{document}
\setlength{\parindent}{0pt}

\def \t {\textrightarrow}
\def \v {\vspace{1em}}
\def \bi {\begin{itemize}}
\def \ei {\end{itemize}}
\def \s[#1] {\section*{#1}}
\def \ss[#1] {\subsection*{#1}}
\def \sss[#1] {\subsubsection*{#1}}

\s[Carducci - traversando la maremma]
Paesaggio interiozzato \t paesaggio che diventa elegia \t non come in D'annunzio in cui il paesaggio è sottoposto a una metamorfosi, in cui la natura viene umanizzata e viceversa \\
Mostrato un viaggio in treno, per cui il mondo lo vede attraverso lo schermo del vetro che fa da limite \\
Sono presenti pilastri degli elementi leopardiani e pascolani \t non è una novita questo incipit con sinestesia, che si configura come molto profonda: dolce paese, ricorda l'universo dell'addio ai monti di manzoni \\
Descrizione del paese molto intima \\
Abitus = tempra ardita, la fierezza, il senso di appartenenza \\
Sdegnoso canto \t è il carducci delle invettive, dei jambi ed epodi \t la sua poesia di solito sdegna qualcosa o qualcuno, pur nella dolcezza del contesto \t e questo riporta a leopardi \\
Odio e amor \t e dice petto, non cuore: descrive l'area della persona che raccoglie le emozioni \\
L'oggetto e il ricordo \t dall'oggetto al ricordo \t montale nei limoni vede questo paesaggio e ricostruisce ? \\
Riconosco non conosco = identificazione di un ricordo, ritorna e ricorda \\
Il sorriso e il pianto sono due opposti, che pero non si escludono \t verso 7 anastrofe e enj. \\
Come se il tempo raggiungesse il tempo zero, la morte e quindi l'eternità \t le tre dimensioni del tempo sono passato, presente e futuro (visione agostiniana) \\
Dejauve forte \t quello che sognai invano è petrarchesco ??? \t lettura petrarchesca molto presente \\
Le colline mettono pace al cuore e sono mute \t velo di ossimoro \t il paesaggio nel cuore e il cuore nel paesaggio \t Palazzeschi, nebbie fumanti \\
Incompiutezza del verso 3 della prima strofa e 2 della seconda, e 10 della terza strofa \t celebrazione compatta del ricordo e del paesaggio, e una rivisitazione sapiente dei luogi dell'anima

\ss[Pianto antico]
Dolore irrecuperabile e morte innaturale \t la morte di un figlio vissuta dal padre \\
Ne l'amore serve per riportare in vita \t anastrofe \\
Il melograno non puo sbocciare \t citazione di A silvia, in cui la sua femminilita non è sbocciata \t anche qua, l'infanzia non viene vissuta \\
Orto solingo \t ortus concluso \t limite, muto e solingo = pleonasmo \\
La natura si riprende, rinasce \\
Fiore della mia pianta \t la pianta sono le generazioni precedenti \t che pero è stat percossa, perche la morte è come un colpo di falce \t come il fieno della sera fiesolana di d'annunzio \\
È inaridita, perchè è privata del senso della vita \t è il sangue dentro il genitore che resta e il figlio che non c'è piu \\
Estremo unico fior \t anche questo pleonasmo \\
Sei sei \t anafora, terra nera perchè arida \\
Lirica suggella il percorso di un carducci umano \\
Si coglie la dimensione personale, intima
\end{document}
